\documentclass[aspectratio=169]{beamer}

% ---------- ENCODING & FONTS ----------
\usepackage[utf8]{inputenc}
\usepackage[T1]{fontenc}

% ---------- THEME & LOOK ----------
\usetheme{Madrid}
\usecolortheme{seahorse}
\usefonttheme{professionalfonts}
\setbeamertemplate{navigation symbols}{}
\setbeamercovered{transparent}
\setbeamertemplate{itemize items}[circle]
\setbeamertemplate{enumerate items}[default]
\setlength{\parskip}{0.25em}

% ---------- COLORS ----------
\usepackage{xcolor}
\definecolor{accent}{RGB}{0,92,184}
\definecolor{soft}{RGB}{233,244,255}
\setbeamercolor{structure}{fg=accent}
\setbeamercolor{title}{fg=white,bg=accent}
\setbeamercolor{frametitle}{fg=accent}
\setbeamercolor{block title}{bg=accent!12,fg=accent}
\setbeamercolor{block body}{bg=soft}

% ---------- PACKAGES ----------
\usepackage{booktabs}
\usepackage{amsmath, amssymb}
\usepackage{array}
\usepackage{multirow}
\usepackage{tabularx}
\usepackage{hyperref}
\hypersetup{colorlinks=true,linkcolor=accent,urlcolor=accent}

% ---------- SAFE FALLBACK FOR TRANSITIONS ----------
\makeatletter
\@ifundefined{transfade}{\newcommand{\transfade}[1][]{}}{}
\makeatother

% ---------- PLACEHOLDER BOX (compile-safe; no images required) ----------
\newcommand{\placeholderbox}[2][0.82\linewidth]{%
  \begin{center}
    \fbox{\begin{minipage}[c][0.46\textheight][c]{#1}
      \centering \vspace{0.25em}
      \textbf{#2}\\[0.6em]
      \emph{Insert figure/diagram here later.}
    \end{minipage}}
  \end{center}
}

% ---------- TITLE ----------
\title{\textbf{Unions and Collective Bargaining \vspace{.5cm}\\  }}
\subtitle{ {\large ECON 300 - Labour Economics}}
%\institute{UNBC}
\author{Muhebullah Karimzada}

\date{Winter 2026}



\begin{document}

% ============================================================
% TITLE
% ============================================================
\begin{frame}[plain]
  \titlepage
\end{frame}

% ============================================================
\section{Learning Objectives \& Main Questions}
% ============================================================

\begin{frame}{Learning Objectives}
\transfade
\begin{itemize}[<+->]
  \item Explain the nature of unions and the role they play in the labour market.
  \item Summarize key measures of union influence in Canada.
  \item Describe the legal framework of unions and collective bargaining and its evolution.
  \item Use economic models to explain why unionization varies across workers, industries, and regions.
  \item Identify union objectives and constraints; apply wage--employment determination models.
\end{itemize}
\end{frame}

\begin{frame}{Main Questions (1/2)}
\transfade
\begin{itemize}[<+->]
  \item What fraction of Canadian workers are union members? How has this changed? How does it compare internationally?
  \item Which workers are most likely to be covered by union contracts?
  \item What determines the overall level of unionization in a labour market?
  \item What motives underlie union behaviour (wages vs.\ job security)?
\end{itemize}
\end{frame}

\begin{frame}{Main Questions (2/2)}
\transfade
\begin{itemize}[<+->]
  \item Can union behaviour be represented by economic models?
  \item How do firms and unions interact in setting employment and wages?
  \item Can theory explain practices like featherbedding and other constraints?
  \item How can we incorporate bargaining power in firm--union interaction?
\end{itemize}
\end{frame}

% ============================================================
\section{Unions: Roles and Types}
% ============================================================

\begin{frame}{What Are Unions?}
\transfade
\begin{itemize}[<+->]
  \item Collective organizations aiming to improve members' well-being via \textbf{collective bargaining}.
  \item Outcomes: wages, benefits, and employment relations.
  \item Additional roles: social and political influence, including financial support to parties.
\end{itemize}
\end{frame}

\begin{frame}{Types of Unions}
\transfade
\begin{itemize}[<+->]
  \item \textbf{Craft unions}: represent a trade/occupation (e.g., printing, construction).
  \item \textbf{Industrial unions}: represent all workers in an industry (e.g., auto, steel, forest).
\end{itemize}
\end{frame}

% ============================================================
\section{Unionization in Canada: Levels and Trends}
% ============================================================

\begin{frame}{Table 14.1: Union Membership and Union Density in Canada, 1920--2019}
\transfade
\scriptsize
\renewcommand{\arraystretch}{1.05}
\begin{tabular}{lccc}
\toprule
\textbf{Year} & \textbf{Union Membership (000s)} & \textbf{Union Membership as a Percentage of Civilian Labour Force} & \textbf{Union Membership as a Percentage of Non-agricultural Paid Workers} \\
\midrule
1920 & 374 & 9.4 & 16.0 \\
1925 & 271 & 7.6 & 14.4 \\
1930 & 322 & 7.9 & 13.9 \\
1935 & 281 & 6.4 & 14.5 \\
1940 & 362 & 7.9 & 16.3 \\
1945 & 711 & 15.7 & 24.2 \\
1951* & 1{,}029 & 19.7 & 28.4 \\
1955 & 1{,}268 & 23.6 & 33.7 \\
1960 & 1{,}459 & 23.5 & 32.3 \\
1965 & 1{,}589 & 23.2 & 29.7 \\
1970 & 2{,}173 & 27.2 & 33.6 \\
1975 & 2{,}884 & 29.9 & 35.6 \\
1980 & 3{,}397 & 29.2 & 35.7 \\
1985 & 3{,}666 & 28.3 & 36.4 \\
1990 & 4{,}031 & 28.5 & 34.5 \\
1995 & 4{,}003 & 27.0 & 34.3 \\
2000 & 4{,}013 & 25.3 & 32.6 \\
2005 & 4{,}369 & 25.2 & 32.3 \\
2010 & 4{,}487 & 24.3 & 31.6 \\
2015 & 4{,}658 & 24.1 & 30.8 \\
2019 & 4{,}892 & 24.2 & 30.4 \\
\bottomrule
\end{tabular}
\end{frame}

\begin{frame}{Figure 14.1: Union Density in Canada and the United States, 1920--2019 — \textit{Placeholder}}
\transfade
\placeholderbox{Figure 14.1 — Union density in Canada vs.\ the United States, 1920--2019.}
\end{frame}

% ============================================================
\section{Legal Framework}
% ============================================================

\begin{frame}{Evolution of the Legal Framework in Canada}
\transfade
\begin{itemize}[<+->]
  \item Pre-Confederation: law discouraged collective bargaining.
  \item 1870s: legal neutrality.
  \item 1944: law encouraged the spread of collective bargaining.
\end{itemize}
\end{frame}

\begin{frame}{Key Features of Canadian Labour Relations Policy}
\transfade
\begin{itemize}[<+->]
  \item Right to join and form unions.
  \item Protection via unfair labour practices legislation.
  \item Established representation system; union as exclusive bargaining agent.
  \item Rights enforced by labour relations boards.
\end{itemize}
\end{frame}

% ============================================================
\section{Union Growth and Incidence}
% ============================================================

\begin{frame}{Measures of Unionization}
\transfade
\begin{itemize}[<+->]
  \item \textbf{Union density}: union members as a percentage of paid workers (\(\approx 34\%\) in Canada).
  \item \textbf{Collective agreement coverage}: percentage of paid workers covered by a collective agreement (\(\approx 34\%\) in Canada).
\end{itemize}
\end{frame}

\begin{frame}{Table 14.2: Union Density and Collective Agreement Coverage in Selected OECD Countries}
\transfade
\scriptsize
\renewcommand{\arraystretch}{1.05}
\begin{tabular}{lcccc}
\toprule
 & \multicolumn{2}{c}{\textbf{Union Membership as \% of Paid Workers}} & \multicolumn{2}{c}{\textbf{Collective Agreement Coverage as \% of Paid Workers}} \\
\cmidrule(lr){2-3}\cmidrule(lr){4-5}
\textbf{Country} & \textbf{1980} & \textbf{2016} & \textbf{1980} & \textbf{2016} \\
\midrule
Australia & 50 & 14 & 84 & 60 \\
Canada & 34 & 29 & 37 & 28 \\
Denmark & 79 & 67 & 82 & 82 \\
France & 20 & 8 & 70 & 99 \\
Germany & 35 & 17 & 85 & 56 \\
Italy & 50 & 34 & 80 & 80 \\
Japan & 31 & 17 & 31 & 17 \\
Netherlands & 35 & 17 & 77 & 79 \\
New Zealand & 69 & 18 & 70 & 15 \\
Sweden & 78 & 62 & 88 & 90 \\
United Kingdom & 52 & 24 & 69 & 26 \\
United States & 22 & 10 & 25 & 12 \\
\bottomrule
\end{tabular}
\end{frame}

\begin{frame}{Figure 14.2: \% of Paid Workers Covered by a Union Contract, by Characteristics — \textit{Placeholder}}
\transfade
\placeholderbox{Figure 14.2 — Coverage by industry, occupation, region, and worker traits.}
\end{frame}

% ============================================================
\section{Who Gets Unionized? Evidence and Drivers}
% ============================================================

\begin{frame}{Patterns of Unionization in Canada}
\transfade
\begin{itemize}[<+->]
  \item Highest in nursing and blue-collar occupations.
  \item Higher in public sector than private sector.
  \item Full-time workers and those above 40 are more likely to be unionized.
\end{itemize}
\end{frame}

\begin{frame}{Demand and Supply of Union Representation}
\transfade
\begin{itemize}[<+->]
  \item \textbf{Demand} depends on expected benefits vs.\ costs.
  \item \textbf{Supply} depends on union organizing and contract administration activities.
\end{itemize}
\end{frame}

\begin{frame}{Benefits and Costs to Workers}
\transfade
\begin{columns}[T,onlytextwidth]
\column{0.5\textwidth}
\textbf{Benefits}
\begin{itemize}[<+->]
  \item Higher wages and nonwage benefits
  \item Greater employment security
  \item Protection against arbitrary management
\end{itemize}
\column{0.5\textwidth}
\textbf{Costs}
\begin{itemize}[<+->]
  \item Union dues and time
  \item Potential income loss (e.g., during strikes)
\end{itemize}
\end{columns}
\end{frame}

\begin{frame}{Supply-Side Considerations}
\transfade
\begin{itemize}[<+->]
  \item Contract administration is costly.
  \item Resources reallocated to higher-return organizing.
  \item Organizing effectiveness varies across workplaces.
\end{itemize}
\end{frame}

\begin{frame}{Determinants Highlighted by Empirical Evidence}
\transfade
\begin{itemize}[<+->]
  \item Canada--U.S. unionization differences largely reflect \textbf{supply} differences.
  \item Similar overall desire for unionization across the two countries.
  \item Influences: social attitudes, legislation, broader economic conditions.
\end{itemize}
\end{frame}

\begin{frame}{Worked Example: Decomposing Changes in Unionization (1/2)}
\transfade
\small
Union density in base period \(= .32 \times 43\% + .68 \times 26\% = 31.4\%\).\\
Union density in final period \(= .24 \times 38\% + .76 \times 32\% = 33.4\%\).\\
Change \(= +2\) percentage points.
\medskip

\scriptsize
\renewcommand{\arraystretch}{1.05}
\begin{tabular}{lcccc}
\toprule
 & \multicolumn{2}{c}{\textbf{Percentage of Employment of Paid Workers}} & \multicolumn{2}{c}{\textbf{Union Density}} \\
\cmidrule(lr){2-3}\cmidrule(lr){4-5}
 & \textbf{Base Period} & \textbf{Final Period} & \textbf{Base Period} & \textbf{Final Period} \\
\midrule
Goods-producing industries & 32 & 24 & 43 & 38 \\
Services-producing industries & 68 & 76 & 26 & 32 \\
All industries & 100 & 100 & 31 & 33 \\
\bottomrule
\end{tabular}
\end{frame}

\begin{frame}{Worked Example: Decomposing Changes in Unionization (2/2)}
\transfade
\small
Counterfactual union density \(= .24 \times 43\% + .76 \times 26\% = 30.1\%\).\\
Change due to structural composition \(= 30.1 - 31.4 = -1.3\) percentage points.
\medskip

\scriptsize
\renewcommand{\arraystretch}{1.05}
\begin{tabular}{lcccc}
\toprule
 & \multicolumn{2}{c}{\textbf{Percentage of Employment of Paid Workers}} & \multicolumn{2}{c}{\textbf{Union Density}} \\
\cmidrule(lr){2-3}\cmidrule(lr){4-5}
 & \textbf{Base Period} & \textbf{Final Period} & \textbf{Base Period} & \textbf{Final Period} \\
\midrule
Goods-producing industries & 32 & 24 & 43 & 38 \\
Services-producing industries & 68 & 76 & 26 & 32 \\
All industries & 100 & 100 & 31 & 33 \\
\bottomrule
\end{tabular}
\end{frame}

\begin{frame}{Industry and Enterprise Characteristics}
\transfade
\begin{itemize}[<+->]
  \item Unionization higher in larger firms, concentrated industries, capital-intensive production, and hazardous jobs.
  \item Causes for higher concentration in large firms: weaker individual action, lower per-worker organizing costs, larger potential gains.
\end{itemize}
\end{frame}

\begin{frame}{Personal Characteristics}
\transfade
\begin{itemize}[<+->]
  \item Less union formation among part-timers.
  \item Women are less likely to be represented by a union.
  \item Higher union interest in blue-collar jobs; unionization tends to decrease with age; low-paid workers show higher interest.
\end{itemize}
\end{frame}

% ============================================================
\section{Union Objectives, Constraints, and Bargaining}
% ============================================================

\begin{frame}{Union Objectives \& Preferences}
\transfade
\begin{itemize}[<+->]
  \item Utility increases with wage and employment; indifference curves are downward sloping and convex (diminishing MRS).
  \item Heterogeneous member preferences; leadership choices and information matter.
\end{itemize}
\end{frame}

\begin{frame}{Union Constraints: Labour Demand}
\transfade
\begin{itemize}[<+->]
  \item Wages negotiated; \textbf{employment chosen by firm} along labour demand.
  \item Equilibrium where an iso-utility curve is tangent to the labour demand curve.
\end{itemize}
\end{frame}

\begin{frame}{Figure 14.3: Union Objectives and Constraints — \textit{Placeholder}}
\transfade
\placeholderbox{Figure 14.3 — Iso-utility curves, labour demand, and tangency at equilibrium.}
\end{frame}

\begin{frame}{Special Objective Cases}
\transfade
\begin{itemize}[<+->]
  \item Maximize wage rate: horizontal indifference.
  \item Maximize employment \(E\): vertical indifference.
  \item Maximize real wage bill: rectangular hyperbolas (weights on wage and employment).
  \item Maximize real economic rent: rectangular hyperbolas from origin.
\end{itemize}
\end{frame}

\begin{frame}{Figure 14.4: Union Objectives -- Special Cases — \textit{Placeholder}}
\transfade
\placeholderbox{Figure 14.4 — Alternative objective functions.}
\end{frame}

\begin{frame}{Firm Isoprofit Curves and Bargaining Range}
\transfade
\begin{itemize}[<+->]
  \item Isoprofit curves: wage--employment combinations yielding equal profits; lower curves imply higher profits.
  \item Preferred outcomes define a bargaining range \([W_a, W_u]\).
\end{itemize}
\end{frame}

\begin{frame}{Figure 14.5: Firm Isoprofit Curves — \textit{Placeholder}}
\transfade
\placeholderbox{Figure 14.5 — Isoprofit map in wage--employment space.}
\end{frame}

\begin{frame}{Figure 14.6: Preferred Outcomes \& Bargaining Range — \textit{Placeholder}}
\transfade
\placeholderbox{Figure 14.6 — Firm vs.\ union preferences and feasible bargaining range.}
\end{frame}

\begin{frame}{Relaxing the Demand Constraint}
\transfade
\begin{itemize}[<+->]
  \item Make labour demand more inelastic.
  \item Influence product market (quotas, tariffs), regulations, minimum wages, full-employment policies.
  \item Limit immigration.
\end{itemize}
\end{frame}

\begin{frame}{Efficient Wage--Employment Contracts}
\transfade
\begin{itemize}[<+->]
  \item Negotiating both wage and employment can be Pareto-improving.
  \item Contract curve \(CC'\): tangencies between union indifference and firm isoprofits; lies to the right of labour demand.
\end{itemize}
\end{frame}

\begin{frame}{Figure 14.7: Efficient vs.\ Inefficient Contracts — \textit{Placeholder}}
\transfade
\placeholderbox{Figure 14.7 — Efficient and inefficient wage--employment contracts.}
\end{frame}

\begin{frame}{Obstacles to Efficiency}
\transfade
\begin{itemize}[<+->]
  \item Information constraints, enforceability issues, shocks to demand/finances, and monitoring costs.
\end{itemize}
\end{frame}

\begin{frame}{Figure 14.8: Inefficient, Approximately Efficient, and Efficient — \textit{Placeholder}}
\transfade
\placeholderbox{Figure 14.8 — Contract efficiency comparison.}
\end{frame}

\begin{frame}{Two Benchmark Models}
\transfade
\begin{itemize}[<+->]
  \item \textbf{Labour demand curve model}: negotiate wage; firm sets employment.
  \item \textbf{Efficient contracts model}: negotiate both wage and employment.
\end{itemize}
\end{frame}

% ============================================================
\section{Union Bargaining Power}
% ============================================================

\begin{frame}{Bargaining Power: Wage-Only Negotiations}
\transfade
\begin{itemize}[<+->]
  \item Power linked to labour-demand elasticity: more inelastic \(\Rightarrow\) greater ability to raise wages.
  \item Power also depends on costs of disagreement for both parties.
\end{itemize}
\end{frame}

\begin{frame}{Figure 14.9: Alternative Aspects of Union Power — \textit{Placeholder}}
\transfade
\placeholderbox{Figure 14.9 — Strike threat, demand elasticity, and wage outcomes.}
\end{frame}

\begin{frame}{Union Power and Labour Supply}
\transfade
\begin{itemize}[<+->]
  \item Strike threat is primary bargaining power source.
  \item \textbf{Craft unions} restrict supply by controlling entry/standards.
  \item \textbf{Professional associations} restrict supply via licensing/certification; wages rise but education periods lengthen.
\end{itemize}
\end{frame}

% ============================================================
\section{Chapter Summary}
% ============================================================

\begin{frame}{Summary (1/2)}
\transfade
\begin{itemize}[<+->]
  \item Definitions, union types, and Canadian legal framework.
  \item Measures of unionization (density and coverage) and international comparisons.
  \item Demand--supply view of union representation; empirical patterns by industry, firm size, occupation, and demographics.
\end{itemize}
\end{frame}

\begin{frame}{Summary (2/2)}
\transfade
\begin{itemize}[<+->]
  \item Union objectives and constraints; labour demand vs.\ efficient contract models.
  \item Bargaining range, isoprofits, and the contract curve.
  \item Bargaining power drivers: demand elasticity and strike threat; roles of craft and professional controls on supply.
\end{itemize}
\end{frame}

\begin{frame}
\centering
\Huge{\textbf{Thank You!}} \\

\end{frame}

\end{document}
